\section{System Overview}

\subsection{Main Stakeholders}

Các bên liên quan chính của hệ thống bao gồm:
\begin{itemize}
    \item \textbf{Du khách (Users)}: Người trực tiếp sử dụng ứng dụng để nhận gợi ý và chơi quest.
    \item \textbf{Chính quyền và cơ quan quản lý du lịch}: Các cơ quan kiểm duyệt nội dung, cung cấp thông tin chính thống về di tích.
    \item \textbf{Cộng đồng địa phương}: Những người đóng góp tri thức, thông tin sự kiện ẩm thực đường phố.
\end{itemize}

\subsection{Main Actors (User Types)}

Hệ thống phân chia thành hai actor chính:
\begin{itemize}
    \item \textbf{Anonymous Tourist}: Xem thông tin địa điểm và bản đồ ở mức cơ bản.
    \item \textbf{Registered Explorer}: Lưu lịch sử đi lại, tham gia quest storyline, upload ảnh check-in và nhận điểm thưởng.
\end{itemize}

\subsection{List of Core Features}

Các tính năng cốt lõi được xây dựng trong hệ thống BeVietnam:
\begin{enumerate}
    \item \textbf{Bảng gợi ý cá nhân hóa (Personalized Feed)}: Đề xuất địa điểm phù hợp theo thời tiết, khoảng cách và sở thích văn hóa.
    \item \textbf{Bản đồ bóng động (Dynamic Bubble Map)}: Trực quan hóa mức độ phù hợp của các điểm đến bằng kích thước vòng tròn trên bản đồ.
    \item \textbf{Tuyến nhiệm vụ cốt truyện (Storyline Quests)}: Giao nhiệm vụ thực địa cho người dùng khám phá lịch sử, di tích.
    \item \textbf{Xác minh ảnh chụp (Capture Judge)}: Sử dụng Vision-Language Model để duyệt ảnh check-in tự động.
\end{enumerate}

\subsection{Explanation of Core Features}

(Chi tiết các tính năng đã được trình bày cụ thể thông qua bảng Input/Output ở Phần 2).

\subsection{Innovation Highlights}


